\chapter{La société de market making pour compte propre}\label{ch:fm:the-proprietary-market-making-firm}

En juillet 2024, le conseil d'un teneur de marché coté de produits négociés en bourse (ETP) a annoncé qu'il cesserait de verser des dividendes réguliers. L'argent devait
rester dans la société en tant que capital de trading~: le capital que ses prime brokers lui imposent de conserver auprès d'eux avant de pouvoir trader, et qui fixe combien elle
peut trader. Dix-huit mois plus tard, ce capital s'élevait à 1\,044~millions d'euros, soit un tiers de plus qu'un an plus tôt, et la société avait emprunté 200
millions de dollars en sus. Une société de trading pour compte propre n'a ni argent de clients ni frais~; tout ce qu'elle peut trader est ce que ses propriétaires y ont laissé, et
ce qu'elle peut emprunter en contrepartie. Combien laisser dans la société, et combien dépenser chaque année pour rester assez rapide afin de gagner quoi que ce soit, sont les deux décisions
que modélise ce chapitre.

\section{Propriété~: sociétés de personnes, sociétés privées et sociétés cotées}

Une société de trading pour compte propre (livre~1, chapitre~1) appartient aux personnes qui l'ont fondée et à celles qu'elles ont admises depuis~; quelques-unes sont
cotées. La propriété détermine trois choses~: qui reçoit le profit, qui peut retirer du capital et quand, et qui supporte une perte.

\begin{definition}[Société de personnes à responsabilité limitée, capital des associés]\label{def:fm:the-proprietary-market-making-firm:llp}
Une \emph{société de personnes à responsabilité limitée}\index{société de personnes à responsabilité limitée} (LLP) est une personne morale, dotée d'une personnalité juridique distincte de celle de ses associés,
dont les associés se partagent le profit selon un pacte d'associés et dont la responsabilité est limitée à ce qu'ils ont apporté. Le
\emph{capital des associés}\index{capital des associés} est le capital que l'associé a apporté ou laissé dans la société, inscrit sur son compte de capital
et remboursable au départ selon les termes du pacte.
\end{definition}

Les sociétés privées détenues par leurs fondateurs et leur personnel fonctionnent sur le fond de la même manière~: le profit est distribué sous forme de dividendes ou de primes, au lieu
d'allocations, et le capital est constitué de \omterm{def:fm:the-proprietary-market-making-firm:ownfunds}{fonds propres} au lieu de \omterm{def:fm:the-proprietary-market-making-firm:llp}{capital des associés}. Une société cotée ajoute des actionnaires extérieurs qui attendent une politique de dividende et
un cours de l'action, ce qui explique pourquoi ses décisions de mise en réserve des bénéfices sont annoncées, comme dans l'accroche, alors que celles d'une société privée ne le sont pas.

\begin{omfigure}
\begin{tikzpicture}[font=\footnotesize,>=stealth,box/.style={draw,rounded corners,minimum width=2.6cm,minimum height=0.62cm,align=center,
  fill=omDef!10}]
\node[box] (rev) at (0,3.2) {revenu\\{\scriptsize capture $\times$ volume}};
\node[box] (prof) at (3.6,3.2) {profit\\{\scriptsize après coûts et salaires}};
\node[box,fill=omProp!12] (alloc) at (7.4,3.2) {allocation\\{\scriptsize par points}};
\node[box,fill=omThm!10] (draw) at (7.4,1.6) {prélèvements\\{\scriptsize versement}};
\node[box] (cap) at (3.6,1.6) {capital des associés\\{\scriptsize conservé}};
\node[box] (tc) at (3.6,0) {capital de trading\\{\scriptsize auprès des prime brokers}};
\node[box] (vol) at (0,0) {volume\\{\scriptsize capital / marge}};
\node[box,fill=gray!12] (tech) at (0,1.6) {technologie\\{\scriptsize préserve la capture}};
\draw[->] (rev) -- (prof); \draw[->] (prof) -- (alloc); \draw[->] (alloc) -- (draw);
\draw[->] (alloc.south west) -- (cap.north east); \draw[->] (cap) -- (tc); \draw[->] (tc) -- (vol);
\draw[->] (vol) -- (tech) node[midway,left,font=\scriptsize] {}; \draw[->] (tech) -- (rev);
\draw[->,dashed,gray] (prof.south west) -- node[above,sloped,font=\scriptsize] {dépenses} (tech.north east);
\end{tikzpicture}

\omcaption{L'exercice annuel d'une société de personnes de trading. Le profit est réparti par points~; ce que les associés ne prélèvent pas reste en capital, que les prime brokers transforment
en capacité de trading~; les dépenses technologiques préservent la capture, et donc le revenu, de l'érosion. Modèle~: \texttt{firm.partnership}.}\label{fig:fm:the-proprietary-market-making-firm:cycle}
\end{omfigure}

Le pacte d'associés fixe les \emph{points} de chaque associé, c'est-à-dire la part de profit qui lui est allouée, et la fraction de cette allocation qui peut être prélevée. Le
capital d'un associé sortant est en général remboursé sur plusieurs années, de sorte qu'un départ ne vide pas le capital de trading d'un seul coup. Ce délai fait
partie de la planification du capital de la société~: un départ est un retrait de capital programmé à l'avance.

\begin{example}[Une allocation]\label{ex:fm:the-proprietary-market-making-firm:alloc}
Quatre associés détiennent 40, 30, 20 et 10 points et prélèvent 80~\% de leurs allocations. Le profit annuel de 116.25~millions de dollars alloue 46.5~millions de dollars au
premier~; elle prélève 37.2~millions de dollars et 9.3~millions de dollars restent sur son compte de capital. Au total, les associés prélèvent 93.0~millions de dollars et laissent 23.25~millions de dollars,
soit un \omterm{def:fm:the-proprietary-market-making-firm:retention}{taux de mise en réserve} de 20~\%.
\end{example}

\begin{definition}[Taux de mise en réserve]\label{def:fm:the-proprietary-market-making-firm:retention}
Le \emph{taux de mise en réserve}\index{taux de mise en réserve} d'une société est la part du profit annuel qu'elle garde en capital au lieu de la distribuer à ses propriétaires~;
son complément à un est le taux de distribution.
\end{definition}

\section{Capital~: à quoi il sert et combien il est exigé}

Une société pour compte propre détient du capital pour trois parties prenantes~: son régulateur, ses membres compensateurs et ses prime brokers, et elle-même. L'exigence du régulateur
protège le marché de la défaillance de la société~; celles des membres compensateurs et des prime brokers les protègent des pertes de la société avant de pouvoir
la liquider~; le coussin propre de la société lui permet de continuer à trader malgré les pertes qu'elle anticipe.

\begin{definition}[Fonds propres, exigence au titre des frais généraux fixes, exigence au titre des facteurs K]\label{def:fm:the-proprietary-market-making-firm:ownfunds}
Les \emph{fonds propres}\index{fonds propres} d'une entreprise d'investissement sont le capital réglementaire qu'elle impute à ses exigences~: principalement des capitaux propres libérés et
les résultats non distribués. Selon les règles européennes applicables aux entreprises d'investissement, son exigence est la plus élevée de ces montants~: un minimum permanent, une \emph{exigence au titre des frais
généraux fixes}\index{exigence au titre des frais généraux fixes}, un quart des frais généraux fixes de l'exercice précédent, et une \emph{exigence au titre des
facteurs K}\index{exigence au titre des facteurs K}, la somme de charges proportionnelles à des mesures de son activité (pour un dealer agissant pour compte propre~: son flux quotidien
de négociation, et le risque de ses positions nettes ou la marge versée à la compensation).
\end{definition}

\begin{dated}{2026-09}{Les règles de capital des sociétés de trading pour compte propre}\label{dat:fm:the-proprietary-market-making-firm:rules}
\textbf{Union européenne.} Regulation (EU) 2019/2033 (the Investment Firms Regulation), article~11~: \omterm{def:fm:the-proprietary-market-making-firm:ownfunds}{fonds propres} au moins égaux au plus élevé des montants suivants~: l'\omterm{def:fm:the-proprietary-market-making-firm:ownfunds}{exigence au titre des
frais généraux fixes} (article~13~: un quart des frais généraux fixes de l'exercice précédent), le minimum permanent (article~14, le capital initial
de la Directive (EU) 2019/2034, article~9~: 750\,000~euros pour une société négociant pour compte propre) et l'\omterm{def:fm:the-proprietary-market-making-firm:ownfunds}{exigence au titre des facteurs K} (article~15). Pour un dealer
les facteurs K comprennent le K-DTF, 0.1~\% du flux quotidien de négociation au comptant et 0.01~\% sur les dérivés, le flux étant la moyenne quotidienne des
achats et des ventes en valeur absolue sur six mois se terminant trois mois avant le calcul (article~33), et soit le K-NPR, risque de position nette, soit
le K-CMG, la marge versée à un membre compensateur (article~21). Les règles MIFIDPRU du Royaume-Uni suivent la même conception.
\textbf{États-Unis.} SEC Rule 15c3-1~: le capital net d'un broker-dealer est sa valeur nette ajustée, entre autres, des décotes sur ses positions en
titres~; l'endettement global ne peut excéder 1\,500~\% du capital net ou, selon la norme alternative, le capital net doit être au moins égal au plus grand
de 250\,000~dollars et de 2~\% des éléments débiteurs agrégés~; un dealer conserve au moins 100\,000~dollars.
\end{dated}

\begin{definition}[Règle de capital net]\label{def:fm:the-proprietary-market-making-firm:netcap}
La \emph{règle de capital net}\index{règle de capital net} (net capital rule) est l'exigence américaine selon laquelle un broker-dealer maintient son capital net, c'est-à-dire sa valeur nette
diminuée des actifs peu liquides et des décotes sur ses positions, au-dessus d'un plancher fixé par son activité, et par son endettement ou les débits de ses clients.
\end{definition}

La conception européenne rattache le capital réglementaire d'une petite société à sa base de coûts, et non à son risque. La société de ce chapitre a des frais généraux fixes de 95
millions de dollars par an (personnel 60~millions de dollars, technologie 35~millions de dollars) et négocie 10~milliards de dollars par jour de titres au comptant~: son \omterm{def:fm:the-proprietary-market-making-firm:ownfunds}{exigence au titre des frais généraux fixes}
s'élève à 23.75~millions de dollars, son K-DTF à 10~millions de dollars, son minimum permanent à 0.75~million de dollars. L'exigence contraignante est celle des frais généraux fixes
(\cref{fig:fm:the-proprietary-market-making-firm:capital})~: une société qui embauche ou achète des serveurs relève son capital réglementaire avant d'avoir négocié une
action de plus.

\begin{omfigure}
\begin{tikzpicture}
\begin{axis}[width=10.5cm,height=5.2cm,xbar,bar width=10pt,xmode=log,log basis x=10,log origin=infty,xmin=0.3,xmax=600,
  ytick={0,1,2,3,4},yticklabels={minimum permanent,frais généraux fixes (un quart),K-DTF au comptant,K-DTF dérivés,exigence des prime brokers},
  y dir=reverse,ymin=-0.6,ymax=4.6,xlabel={\small millions de dollars (échelle logarithmique)},tick label style={font=\footnotesize},table/col sep=comma,
  nodes near coords,nodes near coords style={font=\scriptsize,anchor=west},point meta=rawx,grid=major,grid style={gray!20}]
\addplot[fill=omDef!55,draw=omDef] table[y=k,x=amount]{figdata/desk/02-the-proprietary-market-making-firm/capital.csv};
\end{axis}
\end{tikzpicture}

\omcaption{La société de ce chapitre (frais généraux fixes de 95~millions de dollars par an, 10~milliards de dollars par jour de négociations au comptant)~: les trois parties de l'exigence européenne en \omterm{def:fm:the-proprietary-market-making-firm:ownfunds}{fonds propres},
le K-DTF si le même flux portait sur des dérivés, et le capital de trading qu'exigent ses prime brokers. L'exigence réglementaire est la plus élevée
des trois premières~; celle des prime brokers vaut huit fois celle-ci. Coefficients~: Regulation (EU) 2019/2033~; paramètres de la société illustratifs. Données~:
\texttt{fm\_partner.regulatory}.}\label{fig:fm:the-proprietary-market-making-firm:capital}
\end{omfigure}

La contrainte qui bride la croissance d'un teneur de marché n'est généralement pas celle du régulateur. Ses prime brokers et ses chambres de compensation fixent leurs propres
exigences, au moyen de modèles internes de décote et de marge, et un teneur de marché doit déposer du capital auprès d'eux avant de pouvoir détenir des positions~; le
teneur de marché coté de l'accroche appelle cela son capital de trading et le planifie chaque jour. Dans le modèle du chapitre, les prime brokers exigent un capital égal à 2~\%
du volume quotidien de la société~: 200~millions de dollars soutiennent 10~milliards de dollars par jour, soit huit fois l'exigence réglementaire.

\begin{dated}{2026-09}{Le capital de trading d'une société}\label{dat:fm:the-proprietary-market-making-firm:flow}
Flow Traders Ltd.\ (rapport annuel 2025)~: en juillet 2024, le conseil d'administration a annoncé un plan d'expansion du capital de trading et suspendu le versement
régulier de dividendes~; aucun acompte sur dividende n'a été versé au titre de 2024 ni de 2025. Le capital de trading net (liquidités nettes auprès des chambres de compensation
et des prime brokers, plus la trésorerie en banque) s'élevait à 1\,043.6~millions d'euros au 31~décembre 2025, contre 774.9~millions d'euros un an plus tôt. En 2025, la société a
contracté un prêt à terme de 200~millions de dollars et une facilité de crédit renouvelable de 75~millions de dollars. Ses prime brokers exigent que la société
«~maintienne certains niveaux minimaux de capital~», calculés au moyen de modèles internes de décote et de marge. Résultat net 2025~: 133.6~millions d'euros~; 635~équivalents temps plein à la clôture de l'exercice.
\end{dated}

\begin{method}[Dimensionner le capital d'une société de trading]\label{met:fm:the-proprietary-market-making-firm:size}
\begin{enumerate}
\item Calculer l'exigence réglementaire~: la plus élevée du minimum permanent, d'un quart des frais généraux fixes de l'année précédente et des facteurs K
des flux et des positions prévus (\cref{lst:fm:the-proprietary-market-making-firm:req}).
\item Calculer l'exigence de chaque prime broker et de chaque chambre de compensation pour les positions prévues, selon leurs propres modèles.
\item Ajouter un coussin pour la pire perte que la société prévoit de supporter sans réduire son activité (chapitres 7 et 8).
\item La société a besoin du plus élevé entre l'exigence réglementaire et la somme des exigences de dépôt, augmentée du coussin~; tout ce qui excède est du capital inactif.
\end{enumerate}
\end{method}

\omcode{code/firm/partnership/firm_partnership.py}{78}{86}{L'exigence en fonds propres de la société~: la plus élevée de ses trois parties.}\label{lst:fm:the-proprietary-market-making-firm:req}

\section{Réinvestissement et distribution}

Le modèle de \texttt{firm.partnership} (\cref{lst:fm:the-proprietary-market-making-firm:sim}) fait tourner une société année après année~: le capital de trading fixe le volume
jusqu'à ce que le marché puisse le fournir~; la marge captée par unité échangée dépend de sa technologie~; le bénéfice, après les coûts fixes, les dépenses technologiques et la
rémunération variable, se répartit entre distribution et mise en réserve. La société du chapitre part de 200~millions de dollars, capte un point de base à parité technologique, verse
25~\% de son résultat en rémunération variable, et ses marchés peuvent lui fournir au plus 20~milliards de dollars par jour, ce que soutiennent 400~millions de dollars de capital. La première année, elle
échange 10~milliards de dollars par jour, gagne 250~millions de dollars, consacre 35~millions de dollars à la technologie et 60~millions de dollars à tout le reste, et dégage 116.25~millions de dollars.

\omcode{code/firm/partnership/firm_partnership.py}{126}{143}{Une année de la société en partenariat~: volume issu du capital, marge issue de la technologie, bénéfice réparti
entre distribution et mise en réserve.}\label{lst:fm:the-proprietary-market-making-firm:sim}

\begin{proposition}[Mettre en réserve jusqu'à la saturation du marché]\label{prop:fm:the-proprietary-market-making-firm:retain}
Supposons que le volume vaille $\min(K/m,\bar V)$ pour un capital $K$, une marge $m$ par unité de volume quotidien et un marché de taille $\bar V$, et que le profit par unité
de volume dépasse le taux d'actualisation des associés multiplié par $m$. Alors chaque dollar mis en réserve tant que $K<m\bar V$ rapporte davantage qu'il ne rapporterait hors de la société,
et chaque dollar mis en réserve dès que $K\ge m\bar V$ ne rapporte rien. Parmi les politiques qui ne dépendent que du capital, mettre en réserve la totalité jusqu'à ce que $K$ atteigne
$m\bar V$ et distribuer la totalité ensuite maximise la distribution actualisée.
\end{proposition}

\begin{proof}[Démonstration partielle]
En deçà de la capacité, un dollar mis en réserve ajoute $1/m$ au volume quotidien et rapporte donc $\pi/m$ par an, $\pi$ étant le profit par unité de volume quotidien et par an~; si $\pi/m$ dépasse le taux
d'actualisation, différer la distribution augmente sa valeur actuelle. Au-delà de la capacité, le volume ne bouge plus et un dollar mis en réserve
ne rapporte rien, alors qu'un dollar distribué vaut sa valeur nominale immédiatement. La politique qui remplit la capacité au plus vite puis
distribue tout est donc optimale parmi les politiques dépendant du capital~; les coûts fixes et variables du modèle ne modifient pas l'argument, car
ils ne dépendent pas de la mise en réserve. Une démonstration complète sur toutes les politiques relève d'un programme dynamique (livre~4, chapitre~9).
\end{proof}

Aux paramètres du chapitre, le meilleur \omterm{def:fm:the-proprietary-market-making-firm:retention}{taux de mise en réserve} fixe est de 20~\%, qui vaut 939~millions de dollars de distributions actualisées sur dix ans à 10~\%
(\cref{fig:fm:the-proprietary-market-making-firm:retention})~; à 20~\%, la société sature son marché la sixième année, à 50~\% la troisième. Mettre en réserve la totalité
jusqu'à ce que le capital atteigne 400~millions de dollars et distribuer la totalité ensuite vaut 1\,210~millions de dollars~: les associés renoncent à deux années de distributions et
gagnent 271~millions de dollars de valeur.

\begin{omfigure}
\begin{tikzpicture}
\begin{axis}[width=10.5cm,height=5.2cm,xlabel={\small taux de mise en réserve (\%)},ylabel={\small distribution actualisée (millions de dollars)},xmin=0,xmax=100,
  ymin=0,ymax=1350,tick label style={font=\footnotesize},grid=major,grid style={gray!20},table/col sep=comma,scaled ticks=false]
\addplot[omDef,thick,mark=*,mark size=1.2pt] table[x=retention,y=value]{figdata/desk/02-the-proprietary-market-making-firm/retention.csv};
\draw[omThm,dashed,thick] (axis cs:0,1209.9) -- (axis cs:100,1209.9);
\node[font=\scriptsize,omThm,anchor=south east,fill=white,inner sep=1pt] at (axis cs:98,1215) {mettre en réserve jusqu'à saturation, puis tout distribuer};
\end{axis}
\end{tikzpicture}

\omcaption{Dix années de distributions actualisées à 10~\%, face à un \omterm{def:fm:the-proprietary-market-making-firm:retention}{taux de mise en réserve} fixe (en points), et pour la politique consistant à tout mettre en réserve jusqu'à ce que le
capital de la société remplisse son marché, puis à tout distribuer ensuite (trait en pointillés). Modèle~: \texttt{fm\_partner.retention\_curve}, paramètres
illustratifs.}\label{fig:fm:the-proprietary-market-making-firm:retention}
\end{omfigure}

Un départ est une décision de mise en réserve que la société n'a pas choisie. Si l'associé détenant 30 points quitte la société à la fin de la cinquième année et est remboursé sur trois
années, la société (mettant en réserve 20~\%) dispose du capital pour saturer son marché à la fin de la huitième année au lieu de la sixième. Le remboursement échelonné est ce qui rend le départ supportable~; un
pacte d'associés qui laisserait un associé senior retirer la totalité d'un coup donnerait à celui-ci le pouvoir d'empêcher la croissance de la société.

\section{La course technologique}

La marge captée par un teneur de marché dépend de sa vitesse par rapport aux autres (livre~11, chapitre~9)~: celui qui atteint en premier une cotation périmée,
celui dont la cotation se fait cueillir. La vitesse est relative, donc rester sur place, c'est prendre du retard.

\begin{definition}[Course technologique]\label{def:fm:the-proprietary-market-making-firm:treadmill}
La \emph{course technologique}\index{course technologique} (technology treadmill) est la situation d'une société de trading dont le revenu dépend de sa technologie par rapport à
celle de ses concurrents~: parce qu'ils continuent d'investir, elle doit dépenser à leur rythme pour maintenir constante sa marge captée, et dépenser moins abaisse sa marge captée
même si ses propres systèmes ne changent pas.
\end{definition}

\begin{proposition}[La dépense qui maintient la parité]\label{prop:fm:the-proprietary-market-making-firm:parity}
Supposons que la technologie des concurrents croisse au taux $g$ par an et que le stock technologique $T_t$ de la société se déprécie au taux $d$, de sorte que
$T_{t+1}=(1-d)T_t+s_t$ pour une dépense $s_t$. La technologie relative $T_t/F_t$ de la société, avec $F_{t+1}=(1+g)F_t$, reste constante si et seulement si
$s_t=(g+d)\,T_t$, et sa dépense croît alors au taux $g$.
\end{proposition}

\begin{proof}
$T_{t+1}/F_{t+1}=T_t/F_t$ ssi $(1-d)T_t+s_t=(1+g)T_t$, soit $s_t=(g+d)T_t$~; alors $T_{t+1}=(1+g)T_t$ et $s_{t+1}=(1+g)s_t$.
\end{proof}

Pour $g=15~\%$ et $d=20~\%$, la société doit dépenser chaque année 35~\% de son stock technologique, soit 35~millions de dollars la première année et 123~millions de dollars la dixième année. Si elle
dépense 25~\%, sa technologie relative, et avec elle sa marge captée, tombe à 0.44 de la parité à la dixième année
(\cref{fig:fm:the-proprietary-market-making-firm:treadmill}). Si le marché ne croît pas, la \omterm{def:fm:the-proprietary-market-making-firm:treadmill}{course technologique} finit par l'emporter dans les deux cas~: le
profit du dépensier à parité tombe à zéro la vingtième année, lorsque des dépenses croissant de 15~\% par an rattrapent un revenu qui ne peut pas croître~; celui du dépensier frugal s'annule
la dix-huitième année, sous l'effet d'une marge captée en baisse (\cref{fig:fm:the-proprietary-market-making-firm:profit}). Si le marché croît de 10~\% par an, le dépensier à parité reste
rentable jusqu'à la cinquante-sixième année.

\begin{omfigure}
\begin{tikzpicture}
\begin{axis}[width=10.5cm,height=5cm,xlabel={\small année},ylabel={\small technologie relative aux concurrents},xmin=1,xmax=25,ymin=0,ymax=1.1,
  tick label style={font=\footnotesize},grid=major,grid style={gray!20},table/col sep=comma,
  legend style={font=\scriptsize,at={(0.5,-0.3)},anchor=north,/tikz/every even column/.append style={column sep=8pt}},legend columns=2]
\addplot[omDef,thick] table[x=year,y=rel_steady]{figdata/desk/02-the-proprietary-market-making-firm/treadmill.csv};
\addplot[omThm,thick,dashed] table[x=year,y=rel_lean]{figdata/desk/02-the-proprietary-market-making-firm/treadmill.csv};
\legend{dépense de 35~\% du stock (parité),dépense de 25~\%}
\end{axis}
\end{tikzpicture}

\omcaption{Le tapis roulant~: la technologie des concurrents croît de 15~\% par an et celle de la société se déprécie de 20~\%. Une dépense de 35~\% du stock maintient la parité~;
une dépense de 25~\% laisse décroître la technologie relative, et avec elle la capture. Modèle~: \texttt{fm\_partner.treadmill}.}\label{fig:fm:the-proprietary-market-making-firm:treadmill}
\end{omfigure}

\begin{omfigure}
\begin{tikzpicture}
\begin{axis}[width=10.5cm,height=5cm,xlabel={\small année},ylabel={\small profit (millions de \$)},xmin=1,xmax=22,ymin=-350,ymax=300,
  tick label style={font=\footnotesize},grid=major,grid style={gray!20},table/col sep=comma,restrict x to domain=1:22,
  legend style={font=\scriptsize,at={(0.5,-0.3)},anchor=north,/tikz/every even column/.append style={column sep=8pt}},legend columns=2]
\addplot[omDef,thick] table[x=year,y=profit_steady]{figdata/desk/02-the-proprietary-market-making-firm/treadmill.csv};
\addplot[omThm,thick,dashed] table[x=year,y=profit_lean]{figdata/desk/02-the-proprietary-market-making-firm/treadmill.csv};
\draw[gray] (axis cs:1,0) -- (axis cs:22,0);
\legend{dépense de 35~\% du stock (parité),dépense de 25~\%}
\end{axis}
\end{tikzpicture}

\omcaption{Profit des deux sociétés de la \cref{fig:fm:the-proprietary-market-making-firm:treadmill} dans un marché qui ne croît pas (mise en réserve de 20~\%).
La société économe ne gagne davantage que la première année~; les deux profits atteignent zéro, aux années 18 et 20. Modèle~: \texttt{fm\_partner.treadmill}.}\label{fig:fm:the-proprietary-market-making-firm:profit}
\end{omfigure}

\begin{remark}[Ce que montrent les dépôts, et ce qu'ils ne montrent pas]\label{rem:fm:the-proprietary-market-making-firm:filings}
Le teneur de marché électronique coté du chapitre 1 a déclaré des coûts de communication et de données de 209.4~millions de dollars en 2019 et de 249.2~millions de dollars en 2025, soit une hausse de
19~\%, alors que son \omterm{def:fm:the-economics-of-a-trading-firm:ntr}{revenu net de trading} oscillait entre 1.0~milliard et 2.3~milliards de dollars~; ses coûts ont presque triplé entre 2016 et 2019, mais à la suite de deux
acquisitions, et non du seul tapis roulant. Un dépôt mêle le tapis roulant à la croissance, aux acquisitions et à l'inflation. Le tapis roulant est un mécanisme
que la recherche publiée étaye~: Budish, Cramton et Shim décrivent une course aux armements continue pour la vitesse, où les gains sont annulés par la concurrence, et
Baron, Brogaard, Hagströmer et Kirilenko ont montré que les sociétés qui ont amélioré leur vitesse relative gagnaient davantage (livre~11, chapitre~1).
\end{remark}

\section{Tutoriel~: dix ans d'une société de personnes}\label{tut:fm:the-proprietary-market-making-firm:tut}

\begin{tutorial}
\textbf{Objectif.} Choisir la part du profit annuel qu'une société de personnes de trading doit conserver, compte tenu de ses règles de capital et du tapis roulant du modèle.
\textbf{État final~:} les \cref{fig:fm:the-proprietary-market-making-firm:retention,fig:fm:the-proprietary-market-making-firm:profit}.
\begin{tutsteps}
\item \textbf{L'exigence.} \texttt{fm\_partner.regulatory()} calcule les \omterm{def:fm:the-proprietary-market-making-firm:ownfunds}{fonds propres} requis pour la première année à partir des frais généraux fixes (95~millions de dollars) et
du flux quotidien de transactions (10~milliards de dollars)~: 23.75~millions de dollars, fixés par les frais généraux fixes.
\item \textbf{L'allocation.} \texttt{fp.Partnership} avec des membres de 40, 30, 20 et 10 points alloue 116.25~millions de dollars
(l'\cref{ex:fm:the-proprietary-market-making-firm:alloc})~; \texttt{fm\_partner.accounts()} simule dix ans avec un départ en cinquième année.
\item \textbf{La mise en réserve.} \texttt{fm\_partner.retention\_curve()} simule dix ans pour des \omterm{def:fm:the-proprietary-market-making-firm:retention}{taux de mise en réserve} allant de 0 à 100~\% par pas de 5~\% et
actualise les versements à 10~\%~; \texttt{target\_value()} applique la politique de remplissage puis de versement.
\item \textbf{Le tapis roulant.} \texttt{fm\_partner.treadmill()} simule 25 ans, à 35~\% puis à 25~\% du stock technologique.
\end{tutsteps}
\end{tutorial}

\textbf{Pour aller plus loin.} Faire croître le marché (exercice~7)~; imposer aux prime brokers une marge plus stricte lors d'une année volatile et observer comment le
capital cible évolue.

\section{Construction~: le modèle de la société de personnes}\label{bld:fm:the-proprietary-market-making-firm:partnership}

\begin{build}
\bfield{Rôle} Le côté des propriétaires d'une société de trading pour compte propre~: qui possède le capital, ce que les règles exigent de lui, et ce que le fait de le conserver ou de le verser
change pour la société.
\bfield{Interface} \texttt{firm.partnership}~: \texttt{Member}, \texttt{Partnership} avec \texttt{allocate(profit, payout\_ratio)},
\texttt{depart(name, years)}, \texttt{pay\_departures()}, \texttt{total\_capital}~;
\texttt{own\_funds\_requirement(fixed\_overheads, dtf\_cash, dtf\_deriv, npr, \dots)}~; \texttt{steady\_spend\_share(g, d)}, \texttt{capture(rel, c0, eta)}~; \texttt{Params},
\texttt{simulate(params, retention)} avec un taux ou une politique, \texttt{target\_capital(target)}, \texttt{discounted\_payout(sim, rate, terminal\_multiple)}.
\bfield{Règles} Une perte est allouée selon les points et réduit le capital~; rien n'est prélevé sur une perte. Le capital d'un membre sortant est remboursé en
versements égaux. Le volume ne dépasse jamais le capital rapporté à la marge, ni la taille du marché. Les coefficients sont des données d'entrée.
\bfield{Tests d'acceptation} \texttt{code/firm/partnership/tests/}~: allocations selon les points, pertes et remboursements échelonnés~; l'exigence retient la part la
plus élevée~; une dépense de parité maintient la technologie relative à un, et une dépense plus faible la fait décroître~; le capital évolue comme le profit diminué du versement~; la
politique de capital cible ne verse rien tant que la cible n'est pas atteinte.
\bfield{Pour aller plus loin} Une capture stochastique avec des années déficitaires et un arrêt du trading lorsque le capital passe sous l'exigence~; un prêt à terme qui ajoute du
capital de trading moyennant un coût d'intérêt.
\end{build}

\begin{omsources}
\item Règlement (UE) 2019/2033 sur les exigences prudentielles applicables aux entreprises d'investissement (articles 11, 13--15, 21, 33) et directive (UE) 2019/2034,
article 9.
\item US Code of Federal Regulations, 17 CFR 240.15c3-1 (net capital rule).
\item Limited Liability Partnerships Act 2000 (Royaume-Uni), section 1.
\item Flow Traders Ltd., Annual Report 2025.
\item E.~Budish, P.~Cramton et J.~Shim, ``The high-frequency trading arms race: frequent batch auctions as a market design response'',
\emph{Quarterly Journal of Economics} 130(4), 2015.
\end{omsources}

\section{Exercices}

\begin{exercise}[$\star$]\label{exo:fm:the-proprietary-market-making-firm:1}
Un dealer pour compte propre a supporté l'an dernier des frais généraux fixes de 95~millions de dollars, traite 10~milliards de dollars par jour de titres au comptant et ne détient aucune
position overnight. Calculer les trois composantes de son exigence européenne en \omterm{def:fm:the-proprietary-market-making-firm:ownfunds}{fonds propres} et dire laquelle est contraignante.
\end{exercise}

\begin{exercise}[$\star$]\label{exo:fm:the-proprietary-market-making-firm:2}
La même société traite plutôt 10~milliards de dollars par jour de notionnel de dérivés. Quel est son K-DTF, et l'exigence contraignante change-t-elle~?
\end{exercise}

\begin{exercise}[$\star$]\label{exo:fm:the-proprietary-market-making-firm:3}
Allouer un profit de 116.25~millions de dollars à des membres de 40, 30, 20 et 10 points qui retirent 80~\% de leur allocation. Quel montant chacun retire-t-il, et
qu'est-ce qui reste dans la société~?
\end{exercise}

\begin{exercise}[$\star\star$]\label{exo:fm:the-proprietary-market-making-firm:4}
La technologie des concurrents croît de 15~\% par an et celle de la société se déprécie de 20~\%. Quelle part de son stock doit-elle dépenser pour maintenir la parité, et combien
cela représente-t-il la dixième année si elle dépense 35~millions de dollars la première année~?
\end{exercise}

\begin{exercise}[$\star\star$]\label{exo:fm:the-proprietary-market-making-firm:5}
Le marché est saturé à 20~milliards de dollars par jour et la capture d'un point de base, les coûts fixes sont de 60~millions de dollars et la dépense de parité part de 35~millions de dollars
et croît de 15~\% par an. Dans quelle année la dépense annule-t-elle le profit avant rémunération variable~? Le démontrer sous forme close.
\end{exercise}

\begin{exercise}[$\star\star$]\label{exo:fm:the-proprietary-market-making-firm:6}
Dans la \cref{fig:fm:the-proprietary-market-making-firm:retention}, pourquoi la valeur baisse-t-elle pour des \omterm{def:fm:the-proprietary-market-making-firm:retention}{taux de mise en réserve} supérieurs à 20~\%, et pourquoi la courbe à taux fixe
reste-t-elle toujours sous la ligne pointillée~?
\end{exercise}

\begin{exercise}[$\star\star\star$]\label{exo:fm:the-proprietary-market-making-firm:7}
\emph{Programmation.} Relancer le dépensier à parité sur 60 ans avec un marché qui croît de 10~\% par an (\texttt{fm\_partner.growing\_market}). Dans quelle année son profit tombe-t-il
à zéro pour la première fois, et pourquoi tombe-t-il seulement~?
\end{exercise}

\begin{exercise}[$\star\star\star$]\label{exo:fm:the-proprietary-market-making-firm:8}
\emph{Trouver la faille.} «~Nous avons réduit les dépenses technologiques de 35~\% à 25~\% de notre stock l'an dernier, et le profit est passé de 116 millions de dollars à 124 millions de dollars.
Réduisez-les à nouveau.~»
\end{exercise}

\section{Problème~: garder ou distribuer}

\begin{problem}[{Problème du week-end --- garder ou distribuer}]\label{pb:fm:the-proprietary-market-making-firm:1}
Quatre associés dirigent une société de market making dotée de 200 millions de dollars de capital. Ils doivent décider quelle part du profit de chaque année garder, et quelle part
consacrer à la technologie.

\medskip\noindent\textbf{Partie I --- Les règles.}
\begin{enumerate}
\item Définir les \omterm{def:fm:the-proprietary-market-making-firm:ownfunds}{fonds propres} et énoncer l'exigence européenne au titre des \omterm{def:fm:the-proprietary-market-making-firm:ownfunds}{fonds propres}.
\item Calculer l'\omterm{def:fm:the-proprietary-market-making-firm:ownfunds}{exigence au titre des frais généraux fixes} de la société la première année.
\item Calculer son K-DTF sur 10 milliards de dollars par jour de transactions au comptant, puis sur le même notionnel de dérivés.
\item Quelle exigence est contraignante, et comment se compare-t-elle aux 200 millions de dollars exigés par ses prime brokers~?
\item Énoncer en une phrase le mécanisme de la \omterm{def:fm:the-proprietary-market-making-firm:netcap}{règle de capital net} américaine.
\end{enumerate}

\medskip\noindent\textbf{Partie II --- La société de personnes.}
\begin{enumerate}[resume]
\item Définir une \omterm{def:fm:the-proprietary-market-making-firm:llp}{société de personnes à responsabilité limitée} (LLP) et le \omterm{def:fm:the-proprietary-market-making-firm:llp}{capital des associés}.
\item Calculer le volume, le revenu, les dépenses technologiques et le profit de la première année.
\item Répartir le profit de la première année entre l'associée de 40 points, au taux de distribution de 80~\%~: son prélèvement et ce qu'elle laisse dans la société.
\item L'associé de 30 points quitte la société à la fin de la cinquième année, et il est remboursé sur trois ans. En quelle année la société remplit-elle son marché (\omterm{def:fm:the-proprietary-market-making-firm:retention}{taux de mise en réserve} de 20~\%),
contre quelle année sans ce départ~?
\end{enumerate}

\medskip\noindent\textbf{Partie III --- La mise en réserve.}
\begin{enumerate}[resume]
\item Définir le \omterm{def:fm:the-proprietary-market-making-firm:retention}{taux de mise en réserve}.
\item Donner la distribution actualisée sur dix ans pour des \omterm{def:fm:the-proprietary-market-making-firm:retention}{taux de mise en réserve} de 0, 20, 50 et 80~\%.
\item Donner le meilleur \omterm{def:fm:the-proprietary-market-making-firm:retention}{taux de mise en réserve} fixe et sa valeur.
\item Donner la valeur de la mise en réserve intégrale jusqu'à ce que le marché soit saturé, et expliquer le gain grâce à la \cref{prop:fm:the-proprietary-market-making-firm:retain}.
\item En quelle année le marché est-il saturé pour des \omterm{def:fm:the-proprietary-market-making-firm:retention}{taux de mise en réserve} de 20, 50 et 80~\%~?
\end{enumerate}

\medskip\noindent\textbf{Partie IV --- La \omterm{def:fm:the-proprietary-market-making-firm:treadmill}{course technologique}.}
\begin{enumerate}[resume]
\item Définir la \omterm{def:fm:the-proprietary-market-making-firm:treadmill}{course technologique} (technology treadmill).
\item Établir la part des dépenses qui maintient la parité.
\item Quelle est la capture du dépensier frugal après dix ans, en fraction de la parité~?
\item En quelle année le profit de chaque société atteint-il zéro dans un marché qui ne croît pas~?
\item Énoncer le \emph{résultat nommé}~: la part des dépenses à la parité et son montant à la dixième année, la meilleure politique de mise en réserve et sa valeur face au meilleur
taux fixe.
\item En deux phrases, que doivent faire les associés en premier~?
\end{enumerate}
\end{problem}

\section{Questions d'entretien}

\begin{interviewq}[$\star$ \iqroles{trader, risque}]\label{iq:fm:the-proprietary-market-making-firm:1}
Selon les règles européennes, pourquoi le capital réglementaire d'une petite société de trading pour compte propre peut-il augmenter lorsqu'elle embauche, même si elle ne traite pas davantage~?
\end{interviewq}

\begin{interviewq}[$\star$ \iqroles{trader}]\label{iq:fm:the-proprietary-market-making-firm:2}
Quelle exigence de capital contraint en général la croissance d'un teneur de marché~: celle du régulateur ou celle de ses prime brokers~? Pourquoi~?
\end{interviewq}

\begin{interviewq}[$\star\star$ \iqroles{trader, chercheur}]\label{iq:fm:the-proprietary-market-making-firm:3}
Votre société peut utiliser 400 millions de dollars de capital de trading et en possède 250 millions. Les associés doivent-ils prélever le profit de cette année~?
\end{interviewq}

\begin{interviewq}[$\star\star$ \iqroles{développeur}]\label{iq:fm:the-proprietary-market-making-firm:4}
La technologie des concurrents s'améliore de 10~\% par an, et la vôtre se déprécie de 25~\%. Quelle part de votre stock technologique devez-vous dépenser chaque année pour suivre le rythme~?
\end{interviewq}

\begin{interviewq}[$\star\star$ \iqroles{risque}]\label{iq:fm:the-proprietary-market-making-firm:5}
Pourquoi les conventions d'associés remboursent-elles le capital d'un associé sortant sur plusieurs années~?
\end{interviewq}

\begin{interviewq}[$\star\star\star$ \iqroles{trader, développeur}]\label{iq:fm:the-proprietary-market-making-firm:6}
Une société a réduit son budget technologique, et son profit a augmenté l'année suivante. Que mesureriez-vous avant de conclure que la réduction était justifiée~?
\end{interviewq}
